% =============================================================================
% Module 17: Solutions to Exercises
% =============================================================================
% This file is conditionally included from main.tex via \ifsolutions
% Attestation: ../Solutions/17_Central_Mass_Bias/problems_17.wl \annot{check ids
% quoted below exist in that script}.
% =============================================================================

\section{Solutions to Exercises}
\label{sec:cb_solutions}

\noindent
The Mathematica script
\solnlink{17_Central_Mass_Bias/problems_17.wl}{problems\_17.wl}
verifies every symbolic and numeric answer below.

% ----- Exercise 17.1 -----
\subsection*{Exercise~\ref{ex:cb_two_channels}: The Two Channels}

(a) Take the logarithm of $\Brat=\etamea^{\gp-1}$ and form the
total differential in the two fitted quantities $(\etamea,\gp)$:
\begin{equation*}
  \ln\Brat=(\gp-1)\ln\etamea
  \;\Longrightarrow\;
  \frac{\delta\Brat}{\Brat}
  =\pderiv{\ln\Brat}{\etamea}\,\delta\etamea+\pderiv{\ln\Brat}{\gp}\,\delta\gp
  =(\gp-1)\,\frac{\delta\etamea}{\etamea}+\ln\etamea\;\delta\gp .
\end{equation*}

(b) With $\gp=1.6$:
\begin{align*}
  \text{ratio channel:}\quad &(\gp-1)\,\frac{\delta\etamea}{\etamea}=(0.6)(0.005)=+0.300\%,\\
  \text{slope-drag channel:}\quad &\ln\etamea\,\delta\gp=\ln(0.79)\times0.03=(-0.2357)(0.03)=-0.707\%.
\end{align*}
The channels have opposite signs. The slope channel is larger by a
factor $2.357$ and sets the sign of the net bias,
$\delta\Brat/\Brat=+0.300\%-0.707\%=-0.407\%$.

(c) At fixed fractional shift $\delta\etamea/\etamea$ the ratio
channel contains no $\etamea$ at all. The slope channel contains
$\ln\etamea$, which vanishes as $\etamea\to1$: when the two planes
coincide, a change of slope rescales both Einstein radii identically and
leaves their ratio alone. The two channels of part~(b) have equal
magnitude when $|\ln\etamea|\,\delta\gp=(\gp-1)\,\delta\etamea/\etamea$,
i.e.\ at
$\etamea=\exp[-(\gp-1)(\delta\etamea/\etamea)/\delta\gp]=e^{-0.1}=0.905$.
The slope channel wins for any $\etamea$ further from unity than this,
so it matters most when the two Einstein radii are well separated.
\emph{[Attested: \solnlink{17_Central_Mass_Bias/problems_17.wl}{problems\_17.wl},
checks 1a--1f.]}

% ----- Exercise 17.2 -----
\subsection*{Exercise~\ref{ex:cb_annihilation}: The Annihilation Boundary}

(a) Multiply $\ysrc''=0$ by $\theta^3$:
$(2-\gp)(\gp-1)\theta^{3-\gp}-2u=0$, so
$u=\tfrac12(2-\gp)(\gp-1)\theta^{3-\gp}$. Substituting into $\ysrc'$,
\begin{equation*}
  \ysrc'=1-(2-\gp)\theta^{1-\gp}+\tfrac12(2-\gp)(\gp-1)\theta^{1-\gp}
  =1-\frac{(2-\gp)(3-\gp)}{2}\,\theta^{1-\gp}.
\end{equation*}
Setting this to zero gives
$\theta_\star^{\gp-1}=(2-\gp)(3-\gp)/2$, which is
eq.~\eqref{eq:cb_thetastar}. Then
$\mBH_{\rm crit}=u=\tfrac12(2-\gp)(\gp-1)\theta_\star^{3-\gp}$, which is
eq.~\eqref{eq:cb_mcrit}. Elimination gives a closed form because, on the
curve $\ysrc''=0$, the black-hole term in $\ysrc'$ has the same power of
$\theta$ as the power-law term.

(b) At $\gp=1.6$: $(2-\gp)(3-\gp)/2=0.28$, so
$\theta_\star=0.28^{1/0.6}=0.1198$ and
$\mBH_{\rm crit}=0.12\times0.1198^{1.4}=0.006155$. A direct numerical
solution of the two equations $\ysrc'=\ysrc''=0$ gives the same values to
25 digits. For $\Brat=0.8$ the near-plane fraction is
$\mBH_{1,\rm crit}=0.006155\times0.8^{-7/3}=0.01036$.

(c) Maximising eq.~\eqref{eq:cb_mcrit} over $1<\gp<2$ gives
$\gp=1.4846$ and $\mBH_{\rm crit}=0.00660$. Measured on the plane where
the radial image forms, no power-law slope keeps its radial critical
curve once the central mass exceeds $0.66\%$ of the power-law mass
inside $\thetaE$. The boundary goes to zero at both ends: as $\gp\to2$
because the isothermal profile has no radial critical curve to begin
with, and as $\gp\to1$ because the profile approaches a uniform sheet
with $\kappab=1$, for which $\lambda_r=0$ everywhere. For $\gp$ slightly
above $1$, the power law's $\lambda_r=1-(2-\gp)(\thetaE/\theta)^{\gp-1}$
is only marginally negative over a wide range inside $\theta_r$: at
$\theta_r/10$ it equals $1-10^{\gp-1}$, which is $-0.023$ at $\gp=1.01$
against $-2.2$ at $\gp=1.5$. A small positive term $u/\theta^2$ then
lifts it above zero. The radial critical curve does not move to
small radii: as $\gp\to1$, $\theta_r\to e^{-1}\thetaE\approx0.37\,\thetaE$
and $\theta_\star\to e^{-3/2}\thetaE\approx0.22\,\thetaE$.
\emph{[Attested: checks 2a--2i.]}

% ----- Exercise 17.3 -----
\subsection*{Exercise~\ref{ex:cb_differential}: Why the Ratio Is Biased}

From eq.~\eqref{eq:cb_m2_m1},
$\mBH_2/\mBH_1=\Brat^{-1}\Brat^{2/(\gp-1)}=\Brat^{(3-\gp)/(\gp-1)}$, so
the fixed-slope ratio shift is
\begin{equation*}
  \frac{\mBH_1-\mBH_2}{\gp-1}=\frac{\mBH_1}{\gp-1}\Big(1-\Brat^{(3-\gp)/(\gp-1)}\Big).
\end{equation*}
For $1<\gp<3$ the exponent is positive, and for $0<\Brat<1$ the power is
below one, so the bracket is strictly positive: the near ring always
moves out more than the far one. The shift vanishes if and only if
$\mBH_2=\mBH_1$. That happens in two limits: as $\Brat\to1$, when the
source planes coincide and the black hole is the same fraction of the
enclosed mass on both; and as $(3-\gp)/(\gp-1)\to0$, i.e.\ $\gp\to3$,
when the Einstein radii scale with $\Brat$ in the same way as
$\thetaBH$, so that the fraction is again the same on both planes.
In either case the black hole shifts both total-mass Einstein radii by
the same fraction, and their ratio cancels it. (A fit need not return
the total-mass radii; see the keybox ``The bias is differential''.) For $\gp=1.6$ and
$\Brat=0.8$, $\mBH_2/\mBH_1=0.594$ and the bracket is $0.406$.
\emph{[Attested: checks 3a--3e.]}

% ----- Exercise 17.4 -----
\subsection*{Exercise~\ref{ex:cb_coherent}: A Coherent Systematic}

(a) If every lens has bias $b$, the sample mean has bias
$\tfrac1N\sum_{i=1}^N b=b$, independent of $N$. Independent statistical
errors add in quadrature, so the error of the mean is
$\sigma_1/\sqrt{N}$.

(b) The systematic dominates once $\sigma_1/\sqrt{N}<b$, i.e.\
for $N>(\sigma_1/b)^2=(0.02/0.005)^2=16$.

(c) In a flat universe
$\Brat=(1-\chi_{\rm d}/\chi_1)/(1-\chi_{\rm d}/\chi_2)$, with
$\chi(z)\propto\int_0^z\dd z'/E(z')$ and
$E^2=\Omegam(1+z)^3+(1-\Omegam)(1+z)^{3(1+w)}$. A centred finite
difference about $w=-1$ gives $S_w=\partial\ln\Brat/\partial w=-0.0330$.
Then $\delta w=(-0.005)/(-0.0330)=+0.151$. A $\Brat$ biased low by half a
percent shifts $w$ by $+0.15$.
\emph{[Attested: checks 4a--4d.]}

% ----- Exercise 17.5 -----
\subsection*{Exercise~\ref{ex:cb_real_system}: How Large Is It for a Real Lens?}

(a) $\mBH_1=10^{9}/(2\times10^{11})=0.005$.

(b) From eq.~\eqref{eq:cb_m1crit} with $\Brat=0.8$:
$\mBH_{1,\rm crit}(1.6)=0.0104$, so $\mBH_1/\mBH_{1,\rm crit}=0.48$ and
the radial image survives. $\mBH_{1,\rm crit}(1.9)=0.0017<\mBH_1$, so at
$\gp=1.9$ the far plane has no radial image. The scales show why: at
$\gp=1.9$, $\theta_r=0.077\,\thetaE$, while
$\thetaBH=\sqrt{0.005}\,\thetaE=0.071\,\thetaE$.

(c) Eq.~\eqref{eq:cb_boe} gives, for $\gp=1.6$,
$\delta\etamea/\etamea=3.38\times10^{-3}$ and
$\delta\Brat/\Brat=2.03\times10^{-3}$; for $\gp=1.9$,
$\delta\etamea/\etamea=1.33\times10^{-3}$ and
$\delta\Brat/\Brat=1.19\times10^{-3}$. With $S_w=-0.0330$, the $\gp=1.9$
estimate gives $\delta w=-0.036$. At $\gp=1.9$ there is no radial image
to drag the slope, so if the slope is pinned by other data the ratio
channel is all that remains. Even then eq.~\eqref{eq:cb_boe} assumes the
fit returns the total-mass radii, which holds only for images close to
the rings (\S\ref{sec:cb_boe}).

(d) The coefficient $-4.26$ of \S\ref{sec:cb_worked} belongs to
one axisymmetric geometry with one particular pair of source positions.
Moving the far source alone changed it by a factor of four. Moreover,
$\mBH_1$ here is about half of $\mBH_{1,\rm crit}$, well past the range
where Figure~\ref{fig:cb_bias_vs_m} shows the linear response holding.
A better estimate requires the actual image configuration (positions,
including how close the radial image is to the critical curve), the
model family actually fitted (ellipticity, shear, any freedom in the
inner profile), which images and arcs enter the likelihood, and an
independent estimate of $\gp$ and of the inner mass. With those one
computes $\mathsf{J}$ and $\mathbf{s}$ for that configuration, or
simulates and refits, as in \S\ref{sec:cb_worked}.
\emph{[Attested: checks 5a--5f.]}

% ----- Exercise 17.6 -----
\subsection*{Exercise~\ref{ex:cb_sis_point_mass}: Isothermal Sphere plus Point Mass}

(a) For $\gp=2$, eq.~\eqref{eq:cb_thetatot} reads
$\theta=\thetaE+\thetaBH^2/\theta$, i.e.\
$\theta^2-\thetaE\theta-\thetaBH^2=0$, with positive root
\begin{equation*}
  \theta_{\mathrm{E,tot}}=\frac{\thetaE+\sqrt{\thetaE^2+4\thetaBH^2}}{2}
  =\thetaE\,\frac{1+\sqrt{1+4\mBH}}{2}
  =\thetaE\big(1+\mBH-\mBH^2+\order{\mBH^3}\big).
\end{equation*}
Eq.~\eqref{eq:cb_thetatot_series} at $\gp=2$ has $c_1=1/(\gp-1)=1$ and
$c_2=(\gp-4)/[2(\gp-1)^2]=-1$, in agreement.

(b) $\lambda_r=1-\deriv{}{\theta}\big(\thetaE+\thetaBH^2/\theta\big)
=1+\thetaBH^2/\theta^2>0$ for all $\theta$, so there is never a radial
critical curve, and hence no radial image. Without a radial image, the
data have nothing near the centre with which to trade slope against the
missing point mass. For near-isothermal lenses the radial image is, in
any case, the first casualty of a central mass
($\mBH_{\rm crit}\to0$ as $\gp\to2$). The slope-drag channel then works
only through whatever other constraint fixes $\gp$.

(c) $\theta_{\mathrm{E,tot}}/\thetaE=(1+\sqrt{1.04})/2=1.009902$,
compared with $1+\mBH-\mBH^2=1.0099$.
\emph{[Attested: checks 6a--6e.]}

% ----- Exercise 17.7 -----
\subsection*{Exercise~\ref{ex:cb_central_excess}: A Cusp Is a Point Mass from Outside}

(a) For $\theta>\theta_b$, $\alpha=\bar\kappa\,\theta=(2/\theta)\int_0^\theta\kappab\,\theta'\,\dd\theta'$.
Split the integral at $\theta_b$ and add and subtract the power law
inside $\theta_b$:
\begin{equation*}
  \alpha(\theta)=\frac{2}{\theta}\int_0^{\theta}\kappab_{\rm PL}\theta'\,\dd\theta'
  +\frac{2}{\theta}\int_0^{\theta_b}\big(\kappab_{\rm in}-\kappab_{\rm PL}\big)\theta'\,\dd\theta'
  =\alpha_{\rm PL}(\theta)+\frac{\mBH_{\rm eff}}{\theta}.
\end{equation*}
The second integral is a constant, so its contribution has exactly the
$1/\theta$ form of a point mass. With
$\kappab_{\rm in}\theta'=\kappab_{\rm PL}(\theta_b)\theta_b^{\gamma_{\rm in}-1}\theta'^{\,2-\gamma_{\rm in}}$,
\begin{equation*}
  \mBH_{\rm eff}=2\,\frac{3-\gp}{2}\,\theta_b^{1-\gp}\,\frac{\theta_b^{2}}{3-\gamma_{\rm in}}
  -2\,\frac{3-\gp}{2}\,\frac{\theta_b^{3-\gp}}{3-\gp}
  =(3-\gp)\,\theta_b^{3-\gp}\Big[\frac{1}{3-\gamma_{\rm in}}-\frac{1}{3-\gp}\Big],
\end{equation*}
which is positive exactly when $\gamma_{\rm in}>\gp$.

(b) For $\gp=1.9$, $\gamma_{\rm in}=2.1$, $\theta_b=0.02$:
$\mBH_{\rm eff}=1.1\times0.02^{1.1}\times(1/0.9-1/1.1)=0.0030$, which is
$2.3$ times $\mBH_{\rm crit}(1.9)=0.0013$. Both numbers are fractions on
the same (far) plane. The comparison is valid because
$\theta_b<\theta_\star=0.040<\theta_r=0.077$: the region where the
radial critical curve would form lies outside $\theta_b$, where the lens
is exactly a power law plus a point mass. Inside $\theta_b$ the slope
$\gamma_{\rm in}>2$ gives $\deriv{\alpha}{\theta}<0$, so $\lambda_r>1$
there too, and the deflector has no radial critical curve at all.
(With $\theta_b=0.1$ both $\theta_\star$ and $\theta_r$ would lie
inside the break, and the point-mass comparison would not apply.)

(c) A modest steepening of the inner $2\%$ of the Einstein
radius acts, on all images outside it, like a central mass of $0.3\%$
of the power-law mass inside $\thetaE$, comparable to a large black
hole, and enough to remove the radial image of a near-isothermal
deflector. What biases the fit is the
rigidity of the power law, its inability to put the right amount
of mass in the centre, and not the black hole as such. A black hole is
one physical source of that mismatch; a cusp or a nuclear cluster is
another. Conversely, a model with freedom in its inner profile can
absorb all of them, provided the images lie outside the region where the
excess sits.
\emph{[Attested: checks 7a--7f.]}
